%!TEX root = ms.tex
% Draft (2026-11-04): theory section for the symmetric PiP score. Not yet \input in ms.tex.
% Numerical checks: paper/checks/sym_properties.py; empirical results: paper/FINDINGS.md sections 8-9.
\section*{The PiP score}

\subsection*{Paper percentiles}
Consider an evaluation year $Y$. For an article $j$ published in year $p_j \le Y$, let $c_j(Y)$ be the citations it has
received through $Y$. Its \emph{percentile} is the share of articles published in the same year with strictly fewer
citations at $Y$,
\begin{equation}
y_j(Y) = \frac{\left|\{\,i : p_i = p_j,\ c_i(Y) < c_j(Y)\,\}\right|}{\left|\{\,i : p_i = p_j\,\}\right|} \in [0, 1).
\end{equation}
Uncited articles have $y_j = 0$. Because every article is compared only with articles of the same publication year,
evaluated at the same time, $y_j$ needs no correction for the age of the paper or for citation inflation.

\subsection*{Peers and standing curves}
An author $a$ starts their career with their first article, in year $s_a$. At evaluation year $Y = s_a + A$ (career
age $A$), the author's \emph{peers} $\mathcal{P}$ are the authors with the same start year who are active at $Y$ (at least
three articles and an $h$-index of at least three, both computed from information available at $Y$). For a quality
threshold $\theta \in [0,1]$, let
\begin{equation}
n_a(\theta) = \left|\{\, j \in a : p_j \le Y,\ y_j(Y) \ge \theta \,\}\right|
\end{equation}
be the number of the author's articles at or above percentile $\theta$; $n_a(0)$ is the author's article count. Within
the peer group, the author's standing at threshold $\theta$ is the share of peers with strictly fewer articles at or above
$\theta$,
\begin{equation}
g_a(\theta) = G_\theta\!\left(n_a(\theta)\right), \qquad
G_\theta(m) = \frac{\left|\{\, b \in \mathcal{P} : n_b(\theta) < m \,\}\right|}{|\mathcal{P}|}.
\end{equation}
The \emph{standing curve} $\theta \mapsto g_a(\theta)$ starts at the author's productivity percentile ($\theta = 0$:
standing on the number of articles) and ends at the author's standing on the number of top articles ($\theta \to 1$).

\subsection*{Definition}
The PiP score with cutoff $\tau \in [0, 1)$ is the average standing over the thresholds at or above $\tau$,
\begin{equation}
\mathrm{PiP}_\tau(a) = \frac{1}{1-\tau} \int_\tau^1 g_a(\theta)\, d\theta .
\end{equation}
It is the area under the standing curve to the right of $\tau$, rescaled to $[0,1]$. We use $\tau = 1/2$ throughout:
only articles in the top half of their publication year earn credit. Lowering $\tau$ shifts weight towards the number of
articles; raising it shifts weight towards the number of highly cited articles (Fig.~\ref{fig:tau_sweep}).

\subsection*{Properties}
The following hold for a fixed peer group $\mathcal{P}$ (the score is relative to that reference set).

\paragraph{Range.} $0 \le \mathrm{PiP}_\tau(a) < 1$, with $\mathrm{PiP}_\tau(a) = 0$ exactly when, at almost every threshold
$\theta \ge \tau$, no peer has fewer articles at or above $\theta$ than $a$.

\paragraph{Monotonicity.} Adding an article, or raising the percentile of an existing article, never lowers
$\mathrm{PiP}_\tau$. An article whose percentile stays below $\tau$ leaves the score unchanged. \emph{Proof:} each change
weakly increases $n_a(\theta)$ for every $\theta$, and leaves it unchanged for $\theta > y_j$; $G_\theta$ is
non-decreasing.

\paragraph{No reward for volume alone.} An author with many articles below the 50th percentile of their year scores zero
at $\tau = 1/2$, however many such articles they publish. The earlier, asymmetric score rewarded such records (below).

\paragraph{Credit for the best work.} A single article in the top $1\%$ raises $n_a(\theta)$ by one at every threshold,
including the highest, where most peers have none. In a simulated cohort (checks/sym\_properties.py), one article at the
99th percentile scores $0.45$, while one article at the 50th percentile scores $0$. The $h$-index cannot distinguish
these two records (both have $h = 1$), and neither can the asymmetric score (both score $0$).

\paragraph{Factorization.} For an author in the peer group, let $F(m) = |\{b \in \mathcal{P}: n_b(0) \le m\}|/|\mathcal{P}|$
be the productivity percentile, which is positive because the author counts among their own peers. Then
\begin{equation}
\mathrm{PiP}_\tau(a) = F\!\left(n_a(0)\right) \cdot Q_\tau(a), \qquad
Q_\tau(a) = \frac{1}{1-\tau} \int_\tau^1 \frac{g_a(\theta)}{F\!\left(n_a(0)\right)}\, d\theta .
\end{equation}
The \emph{quality multiplier} $Q_\tau$ compares the author's standing on counts of good articles with their standing on
the raw count. It exceeds one when the author ranks higher among peers on good articles than on volume: for 30--36\% of active
authors in our data, depending on career age. Taking logarithms splits the variance of $\log \mathrm{PiP}_\tau$ into a
productivity part, a quality part and their covariance. For $\tau = 1/2$, the quality multiplier accounts for about
three quarters of the variance and productivity for about 30\% (the covariance is slightly negative). The asymmetric
score is the reverse, with productivity at 59--77\%.

\paragraph{Consistency.} $\mathrm{PiP}_\tau$ is not consistent in the sense of \cite{waltman2012inconsistency}: adding
the same article to two authors can reverse their order, because $G_\theta$ is not linear. Reversals are rare (0.7\% of
random pairs in the simulation, against 1.3\% for the asymmetric score); the $h$-index has the same property.

\subsection*{Relation to other indicators}
Writing indicators in terms of the count curve $\theta \mapsto n_a(\theta)$ places PiP among existing ones.
\begin{itemize}
  \item \textbf{Top-$x\%$ counts} (e.g., the number of articles in the top 10\%) read the count curve at a single
  threshold, $n_a(1-x)$, without comparing it with peers.
  \item \textbf{I3} \cite{leydesdorff2011integrated} sums the article percentiles,
  $\sum_j y_j = \int_0^1 n_a(\theta)\, d\theta$: the same integral without peer normalization, so it grows without bound
  with the number of articles.
  \item \textbf{The $h$-index} reads one point of the analogous curve in citation space: the largest $k$ such that $k$
  articles have at least $k$ citations, the corner where the count curve meets the diagonal. PiP integrates the whole
  peer-normalized curve.
  \item \textbf{The asymmetric PiP} used in earlier versions of this work compares the author's count at or above
  $\theta$ with peers' \emph{total} counts, $\int_0^1 F(n_a(\theta))\, d\theta = \sum_k y_{(k)}\,[F(k)-F(k-1)]$, where
  $y_{(1)} \ge y_{(2)} \ge \dots$ are the author's percentiles in decreasing order. It factors as $F(n_a(0))$ times a
  weighted mean of the author's own percentiles, but it gives an author's best articles little or no weight whenever
  peers have more articles than the rank of those articles (for the median Google Scholar author, the top four articles
  receive zero weight), and a cutoff $\tau$ applied to it drifts across cohorts.
\end{itemize}

\subsection*{Computation}
We evaluate the integral on the grid $\theta_k = k/40$, $k = 0, \dots, 39$. On the asymmetric score, where the exact
value is available, the grid sum and the exact score correlate at $0.9999$ (mean difference $+0.0075$).
